\chapter{Simular y ajustar}
\label{ch:simular}

Este capítulo cubre el \emph{modelo directo} (simulación) y el \emph{ajuste
discreto}, ilustrados con el caso más sencillo y a la vez más importante: un
calibrado de $\alpha$-Fe. Recuerde la distinción del
capítulo~\ref{ch:intro}: \textbf{simular} produce un espectro a partir de
parámetros que usted fija; \textbf{ajustar} infiere esos parámetros a partir del
espectro medido.

\section{El modelo directo}

Fitbauer construye el espectro como transmisión:
\[
  T(v) \;=\; b + s\,v \;-\; \sum_{k} A_k(v),
\]
donde $b$ es la línea base, $s$ una pendiente opcional y $A_k(v)$ la
\emph{absorción} de cada componente. Las tres formas discretas básicas son:

\begin{center}
\begin{tabular}{@{}lp{0.6\textwidth}@{}}
\toprule
\textbf{Componente} & \textbf{Parámetros físicos} \\
\midrule
\textbf{Singlete} & $\delta$ (desplazamiento isomérico), $\Gamma$ (anchura),
profundidad. Un único pico. \\[3pt]
\textbf{Doblete} & $\delta$, $\dEQ$ (desdoblamiento cuadrupolar), $\Gamma$,
profundidad. Dos líneas. \\[3pt]
\textbf{Sexteto} & $\delta$, $\dEQ$, $\BHF$ (campo hiperfino), anchuras
$\Gamma_{1,2,3}$, profundidad e intensidades relativas. Seis líneas. \\
\bottomrule
\end{tabular}
\end{center}

\autofig[\textwidth]{fig_sim_components}{Simulación de los tres componentes
discretos: singlete (un pico), doblete ($\dEQ = 0{,}80$~mm/s) y sexteto
($\BHF = 33$~T con relación de intensidades $3:2:1$).}

\subsection{Las matemáticas del modelo}
\label{sec:modelo-mates}

Cada componente es una \textbf{suma de líneas} de forma lorentziana. Una línea
centrada en $v_0$ con anchura $\Gamma$ (la \emph{anchura total a media altura},
FWHM —el parámetro que se ajusta—) vale
\[
  L(v;\,v_0,\Gamma) \;=\; \frac{(\Gamma/2)^2}{(v-v_0)^2 + (\Gamma/2)^2},
\]
(máximo $1$ en $v=v_0$; la semianchura es $\Gamma/2$), y la absorción de un
componente es
\[
  A_k(v) \;=\; d_k \sum_{j} w_j \, L(v;\,v_{0,j},\Gamma_j),
\]
con $d_k$ la \emph{profundidad} y $w_j$ las \emph{intensidades} relativas de sus
líneas. Lo que distingue a cada tipo es \textbf{dónde caen las líneas} y
\textbf{cuántas hay}:

\begin{itemize}[nosep]
  \item \textbf{Singlete}: una línea en $v_0 = \delta$. El desplazamiento
        isomérico $\delta$ mueve todo el patrón.
  \item \textbf{Doblete}: dos líneas en $v_0 = \delta \pm \tfrac{1}{2}\dEQ$. El
        cuadrupolo $\dEQ$ es la separación entre ambas.
  \item \textbf{Sexteto}: seis líneas
        \[
          v_{0,j} \;=\; \delta \;+\; \tfrac{1}{2}\,q_j\,\dEQ \;+\;
          p_j\,\frac{\BHF}{33{,}0},
        \]
        donde $p_j = (\pm0{,}839,\ \pm3{,}084,\ \pm5{,}329)$~mm/s son las
        posiciones publicadas de $\alpha$-Fe a 33~T (por eso el patrón
        \emph{escala} con $\BHF/33$) y $q_j=(+1,-1,-1,-1,-1,+1)$ el patrón de
        primer orden del cuadrupolo. Las intensidades siguen la relación
        $3:x:1:1:x:3$ (§\ref{sec:textura}).
\end{itemize}

El espectro completo es la transmisión $T(v) = b + s\,v - \sum_k A_k(v)$. Esta
parametrización —posiciones ligadas a $\delta$, $\dEQ$, $\BHF$ con significado
físico— es la que permite \emph{leer} el material directamente de los parámetros
ajustados. El desarrollo completo (tratamiento del hamiltoniano
magnético+cuadrupolar, órdenes superiores) está en el anexo~\ref{ap:mates}.

\subsection{Perfil de línea y modelo de absorbente}

En el panel de calibración se elige el \menu{Perfil de línea}:

\begin{itemize}[nosep]
  \item \textbf{Lorentziana}: la forma natural de la línea Mössbauer.
  \item \textbf{Voigt}: convolución de Lorentziana con una gaussiana de anchura
        $\sigma$, útil cuando hay un ensanchamiento instrumental o una pequeña
        dispersión no resuelta.
\end{itemize}

El \menu{Absorbente} modela el grosor de la muestra. La diferencia está en cómo
se convierte la absorción $A(v)$ en transmisión:

\begin{itemize}[nosep]
  \item \textbf{Delgado (lineal)}: $T \approx 1 - A(v)$. La absorción es
        proporcional a la cantidad de Fe; válido para muestras poco densas.
  \item \textbf{Grueso (saturación)}: $T \propto e^{-A(v)}$ (forma tipo integral
        de transmisión). Al crecer el grosor, las líneas \textbf{saturan} —dejan
        de crecer linealmente y se ensanchan—; este modo lo corrige para no
        sesgar áreas ni anchuras. La corrección completa está en el
        anexo~\ref{ap:espesor}.
\end{itemize}

\begin{truco}
Ante muestras densas o líneas muy profundas ($>15$--$20\%$), use \textbf{Grueso}:
con el modelo delgado las áreas relativas (y por tanto las proporciones de fase)
saldrían sesgadas.
\end{truco}

\subsection{Intensidades y textura del sexteto}
\label{sec:textura}

Las seis líneas del sexteto siguen la relación de áreas $3:x:1:1:x:3$. El valor
$x$ (líneas 2 y 5) depende de la orientación de los momentos respecto al haz:

\begin{itemize}[nosep]
  \item En modo \textbf{libre}, las intensidades relativas se ajustan.
  \item En modo \textbf{textura}, un único parámetro $t$ fija $x$ de forma física
        (relación $R_{23}$, ángulo $\theta$), útil para muestras orientadas o
        láminas. El panel muestra $\theta$ y $R_{23}$ derivados de $t$.
\end{itemize}

El tratamiento del \textbf{cuadrupolo} en el sexteto puede ser de primer orden
(perturbativo) o completo; se selecciona en \menu{Ajuste > Modo de ajuste}. El
detalle (tratamiento tipo Kündig del hamiltoniano magnético+cuadrupolar) está en
el anexo~\ref{ap:mates}.

\section{Simular sin ajustar}

La \textbf{simulación en vivo} es una de las señas de identidad de Fitbauer: en
lugar de lanzar un ajuste a ciegas, se construye el modelo a mano y se ve el
espectro cambiar \textbf{al instante} con cada parámetro. Para explorar el
modelo:

\begin{enumerate}[nosep]
  \item Elija la forma de cada componente en su panel (\menu{Forma:}) y actívelo.
  \item Fije los parámetros con los controles deslizantes o escribiendo el valor;
        cada control tiene una barra tipo \emph{slider} sincronizada con la caja
        numérica.
  \item El gráfico se actualiza \textbf{en vivo} con el modelo resultante y sus
        subespectros.
\end{enumerate}

\paragraph{Para qué sirve simular.} La simulación no es un adorno; cumple varias
funciones prácticas:

\begin{itemize}[nosep]
  \item \textbf{Sembrar el ajuste}: un ajuste no lineal necesita un punto de
        partida razonable o cae en un mínimo falso. Acercar el modelo «a ojo»
        antes de \menu{Ajustar} hace el ajuste más rápido y fiable.
  \item \textbf{Entender el efecto de cada parámetro}: mover $\delta$ desplaza
        todo el patrón; $\dEQ$ separa las líneas por pares; $\BHF$ abre o cierra
        el sexteto; $\Gamma$ ensancha; la profundidad hunde. Verlo en vivo asienta
        la intuición física.
  \item \textbf{Prever un experimento o comparar hipótesis}: ¿cómo se vería una
        mezcla de dos fases con estas proporciones? Se simula y se compara con un
        espectro real superpuesto (§\ref{sec:comparar} del capítulo anterior).
  \item \textbf{Docencia}: ilustrar de dónde sale cada línea sin necesidad de
        datos.
\end{itemize}

Cada componente se dibuja como un \textbf{subespectro} propio (con su relleno de
área opcional), de modo que se ve la contribución de cada uno al total. El
\textbf{candado} de cada parámetro lo \emph{fija}: útil para simular variando solo
una magnitud, o para que el ajuste posterior no mueva lo que usted ya conoce.

\begin{nota}
Simular y ajustar comparten exactamente el mismo modelo directo: lo que ve al
simular es, línea por línea, lo que el optimizador intentará encajar. Por eso una
buena simulación inicial es la mejor inversión antes de ajustar.
\end{nota}

\guicap{gui_component_panel}{Panel de un componente sexteto con sus controles de
$\delta$, $\dEQ$, $\BHF$, anchuras e intensidades, y los candados de fijado.}

\section{Ajustar un calibrado de $\alpha$-Fe}
\label{sec:ajustar-alphafe}

El $\alpha$-Fe es el caso de referencia. El procedimiento:

\begin{enumerate}[nosep]
  \item Cargue el espectro (\menu{Archivo > Cargar…}).
  \item Active un único componente de forma \textbf{Sexteto}.
  \item Dé valores iniciales aproximados: $\delta \approx 0$, $\dEQ \approx 0$,
        $\BHF \approx 33$~T, $\Gamma \approx 0{,}30$~mm/s.
  \item \menu{Ajuste > Ajustar}.
\end{enumerate}

El resultado debe reproducir las seis líneas con $\BHF \approx 33$~T y un
$\chi^2$ reducido próximo a 1, como en la figura~\ref{fig:fig_calibration_fit}
(capítulo~\ref{ch:calibracion}). Ese ajuste es, precisamente, el que se marca
como calibración (§\ref{sec:definir-calib}).

\begin{truco}
Si quiere que el ajuste refine además la propia \param{Vmax} contra las
posiciones de $\alpha$-Fe, active \emph{Ajustar Vmax}. Para evitar
degeneración, Fitbauer fija automáticamente el $\BHF$ activo cuando se ajusta la
velocidad (no pueden refinarse ambos a la vez).
\end{truco}

\section{Opciones de ajuste}
\label{sec:opciones-ajuste}

Los criterios estadísticos se eligen en el menú \menu{Ajuste} / panel de opciones:

\begin{itemize}[nosep]
  \item \menu{Verosimilitud}: \textbf{Gauss} ($\chi^2$ con la $\sigma$ del dato) o
        \textbf{Poisson} (IRLS, con la $\sigma$ del modelo). Poisson es más
        apropiada con pocas cuentas.
  \item \menu{Pérdida robusta}: \textbf{Lineal} (mínimos cuadrados),
        \textbf{Soft L1} (suave) o \textbf{Huber}, que restan peso a los puntos
        atípicos (\emph{outliers}).
  \item \textbf{Multistart}: repite el ajuste desde varios puntos de partida para
        no quedar atrapado en un mínimo local.
  \item \menu{Propagar $\sigma$ de calibración (vmax)}: propaga la incertidumbre
        de la \param{Vmax} calibrada a los parámetros ajustados.
\end{itemize}

\begin{nota}
Cuando se activa \emph{Ajustar Vmax}, la calibración interviene en el ajuste; por
eso conviene tener claro qué calibración está fijada (capítulo~\ref{ch:calibracion}).
\end{nota}

\subsection{Restricciones y presets físicos}

\menu{Ajuste > Restricciones entre parámetros…} permite \textbf{ligar}
parámetros (por ejemplo, igualar anchuras entre componentes o imponer relaciones
de intensidad). \menu{Ajuste > Presets físicos de restricciones…} aplica
conjuntos de restricciones habituales de un vistazo. Reducir el número de
parámetros libres estabiliza el ajuste y evita degeneraciones.

\section{Calidad del ajuste y errores}
\label{sec:calidad}

Un ajuste no termina cuando el optimizador converge: hay que \textbf{juzgar si es
bueno} y \textbf{cuantificar la incertidumbre}. Fitbauer reúne para ello varios
indicadores en el panel de información y varios métodos de error.

\subsection{Los estadísticos y para qué sirven}

\begin{center}
\begin{tabular}{@{}lp{0.66\textwidth}@{}}
\toprule
\textbf{Indicador} & \textbf{Qué mide / para qué} \\
\midrule
$\chi^2$ & Suma de residuos al cuadrado ponderados por el ruido. Baja al mejorar
el ajuste, pero depende del número de puntos. \\[3pt]
$\chi^2$ \textbf{reducido} & $\chi^2$ por grado de libertad. Es el \textbf{juez
principal}: $\approx 1$ = ajuste compatible con el ruido; $\gg 1$ = el modelo no
reproduce los datos (faltan componentes, mal centro, mala calibración);
$\ll 1$ = sobreajuste o ruido sobreestimado. \\[3pt]
\textbf{dof} (grados de libertad) & Nº de puntos menos parámetros libres. Contexto
para el $\chi^2$. \\[3pt]
\textbf{AIC} y \textbf{BIC} & Criterios de información: premian el buen ajuste y
\textbf{penalizan añadir parámetros}. Sirven para \textbf{comparar modelos}
(p.\,ej.\ ¿dos sitios o tres?): \emph{menor es mejor}. Evitan el autoengaño de que
«más componentes siempre ajustan mejor». \\[3pt]
\textbf{Correlaciones} & Fitbauer avisa de pares de parámetros muy correlacionados
($|r|$ alto): señal de \textbf{degeneración} (el ajuste no puede separarlos) o de
que sobran parámetros. La solución suele ser fijar o ligar alguno
(§\ref{sec:opciones-ajuste}). \\
\bottomrule
\end{tabular}
\end{center}

\begin{truco}
No persiga un $\chi^2_r$ exactamente 1 añadiendo componentes: vigile AIC/BIC y las
correlaciones. Un modelo con $\chi^2_r=1{,}1$ y parámetros bien definidos es mejor
que uno con $\chi^2_r=1{,}0$ lleno de componentes correlacionados.
\end{truco}

\subsection{Estimación de incertidumbres}

Los estadísticos dicen si el ajuste es bueno; los \textbf{errores} dicen cuánto se
puede confiar en cada parámetro. Fitbauer ofrece tres métodos, de menor a mayor
coste y fiabilidad:

\begin{center}
\begin{tabular}{@{}p{0.34\textwidth}p{0.58\textwidth}@{}}
\toprule
\textbf{Método} & \textbf{Qué hace / cuándo usarlo} \\
\midrule
Covarianza (por defecto) & Errores a partir de la matriz de covarianza del ajuste.
Instantáneo; válido si las incertidumbres son aproximadamente gaussianas y las
correlaciones moderadas. Es la primera estimación. \\[4pt]
\menu{Bootstrap errores (MC)…} & Reajusta muchas réplicas del espectro
remuestreadas según el ruido (Monte Carlo) y mide la dispersión de cada
parámetro. Más robusto ante ruido y no linealidad; cuesta más tiempo. \\[4pt]
\menu{Errores por verosimilitud perfilada…} & Recorre cada parámetro reajustando
los demás y ve cuánto empeora el $\chi^2$ (\emph{profile likelihood}). Da
intervalos \textbf{asimétricos} realistas; los más fiables cuando las
incertidumbres \textbf{no} son gaussianas o hay fuertes correlaciones. \\
\bottomrule
\end{tabular}
\end{center}

\begin{nota}
Regla práctica: use la \textbf{covarianza} durante el trabajo exploratorio y
reserve \textbf{bootstrap} o \textbf{verosimilitud perfilada} para las cifras
que vaya a publicar, sobre todo si el panel avisa de correlaciones altas.
\end{nota}

\section{Herramientas de ayuda al ajuste}

\begin{itemize}[nosep]
  \item \menu{Ajuste > Inicializar desde mínimos} y \menu{Autoajustar desde
        mínimos}: detectan las posiciones de línea y proponen un punto de
        partida. El \textbf{editor de mínimos} permite además \emph{curarlos}
        sobre el propio gráfico: clic para añadir un mínimo, clic sobre un
        marcador para activarlo/desactivarlo, y un contador de contribuciones por
        mínimo.
  \item \menu{Ajuste > Buscar centro}: refina el centro de folding.
  \item \menu{Ajuste > Fijar todos} / \menu{Liberar todos}: bloquean o liberan
        todos los parámetros de una vez.
  \item \menu{Ajuste > Ajuste en serie…}: ajusta una tanda de espectros por
        \emph{warm-start} (cada ajuste parte del anterior), ideal para series de
        temperatura o de tratamiento.
  \item \menu{Ajuste > Deshacer ajuste}: revierte al estado previo.
  \item \menu{Ajuste > Historial de ajustes…}: tabla con los últimos ajustes
        terminados (hora, fichero, modo, $\chi^2_r$ y componentes). Cada entrada
        guarda una \textbf{instantánea completa de la sesión}, de modo que
        \param{Restaurar} recupera aquel ajuste íntegro (datos, calibración,
        modelo y opciones). El historial persiste entre sesiones y su tamaño
        máximo es configurable desde el propio diálogo.
  \item \menu{Ajuste > IA local Ollama: sugerir inicio…}: si tiene un modelo
        local, propone valores iniciales a partir de un resumen del espectro.
\end{itemize}

\guicap{gui_minima_editor}{Editor semi-manual de mínimos sobre el gráfico: los
marcadores incluidos (rojos) y excluidos (huecos) se editan con el ratón y con la
lista lateral.}

\section{Componentes de relajación magnética}

Además de singlete/doblete/sexteto, Fitbauer incluye componentes para
\textbf{relajación superparamagnética} y dinámica de espín, útiles en
nanopartículas cerca de la temperatura de bloqueo:

\begin{itemize}[nosep]
  \item \textbf{Relajación} y \textbf{Blume--Tjon}: modelos de línea con una tasa
        de relajación $\nu$ (parámetro $\log_{10}\nu$).
  \item \textbf{Néel (tamaño)}: distribución de tamaños lognormal con anisotropía
        $K_{\text{eff}}$ y temperatura de bloqueo; \menu{Ajuste > Ajuste global
        Néel (multi-T)…} ajusta a la vez varias temperaturas.
\end{itemize}

\begin{nota}
Estos modelos comparten la interfaz de componente (se eligen en \menu{Forma:}).
Su fundamento numérico se desarrolla en el anexo~\ref{ap:relax} (\emph{Modelos de
relajación magnética}).
\end{nota}
